\documentclass[11pt]{article}
\usepackage[margin=1in]{geometry}
\usepackage{amsmath,amssymb}
\usepackage{booktabs}
\usepackage{longtable}
\usepackage{hyperref}
\usepackage{xcolor}
\hypersetup{colorlinks=true, linkcolor=blue, urlcolor=blue}

\title{Discretization Error, Basin Selection, and Pontryagin Residuals\\
       in Direct-Collocation Low-Thrust Orbit Transfers}
\author{M.\ Casey}
\date{Begun 2026--07--27; last revised \today}

\newcommand{\mf}{m_f}
\newcommand{\dV}{\Delta V}

\begin{document}
\maketitle

\begin{abstract}
We set out to answer a question that none of our existing verification
machinery addresses: \emph{is the collocation mesh we solve on fine enough
that the discrete solution is a faithful approximation of the continuous
optimal control problem?} The interior-point solver certifies the nonlinear
program, and our first-order optimality layer certifies the Karush--Kuhn--Tucker
conditions \emph{of that program}, but neither says anything about the gap
between the program and the continuous problem it discretizes. That gap is
discretization error. We measure it by solving the same problem on a ladder of
successively finer meshes and observing the rate at which the answer settles.

The measurement produced a partly unexpected result, and a first draft of this
document over-claimed it in a way that external review caught; see
Section~\ref{sec:corrections}. What we can establish is this. \emph{No} row
yielded an established asymptotic order, because every mesh ladder crossed a
change in the discrete switch structure and an order fitted across such a
change is meaningless. Mesh adequacy therefore remains unestablished. What is
firmly established instead concerns \emph{which local minimum the solver
reaches}: every certified Earth row we examined is beaten \emph{on its own
unmodified mesh}---by $1.0$ to $6.0$~kg, with lower $\Delta V$ as well---by a
solution reached from a different starting guess. That comparison is made at
fixed discretization, so it cannot be a discretization effect. The problem is
non-convex, its basins are reachable from very small changes in the initial
guess, and basin search is demonstrably a higher-yield intervention than mesh
refinement. Whether it \emph{dominates} discretization error quantitatively is
not something this study established.

The investigation then changed method. Because mesh convergence is inherently
\emph{comparative}---it needs two solutions and therefore a proof they are the
same---we replaced it with an \emph{absolute} measurement of a single
solution: does the interpolated trajectory satisfy the continuous Pontryagin
necessary conditions, and to what residual? That measurement succeeded where
the ladder failed, and it produced a sharp answer.
Section~\ref{sec:phase3} reports it: the residual is the local truncation
error of the trapezoidal rule, scaling as $h^3$ with a coefficient fixed by the
orbital dynamics, peaking at apogee, and essentially independent of the
bang-bang switch structure. An attempt to convert that residual into kilograms then failed review for the
second time, and Section~\ref{sec:corrections2} records why. What survives is a
\emph{lower bound}: two $10$~N solutions at identical node count, differing only
in node placement, have final masses $0.1456$~kg apart. A defensible point
estimate of the discretization error remains out of reach, and the reason is
now understood---these solutions are so basin-fragile that a $2\%$ change in the
mesh reorganizes their switch structure, which is itself the study's most
robust result.
\end{abstract}

\tableofcontents

\section{Corrections after external review}
\label{sec:corrections}

The first draft of this document was reviewed by two independent language
models tasked with attacking its \emph{reasoning} rather than its prose. Both
identified the same fatal flaw, independently, and both named it the single
weakest link. It is recorded here rather than quietly fixed, because the error
is instructive: the study broke its own stated rule, and nothing in the code
prevented it.

\paragraph{The error.} Section~\ref{sec:richardson-limits} states that
Richardson extrapolation must not be applied across a branch change. The $1$~N
ladder has switch counts $173, 171, 171, 171$: the first difference spans a
change in the discrete structure. The first sliding window, $p = 1.596$, was
computed across it anyway, and the agreement between that window and the valid
one ($p = 1.228$) was then cited as evidence that an asymptotic regime had been
reached. Only \emph{one} valid window exists, and a single window agrees with
itself trivially.

\paragraph{What was retracted.} Enforcing the rule in code (rather than in a
comment) gives: zero valid windows at $10$~N, one at $2.5$~N reading $-1.67$,
and one at $1$~N reading $1.228$. Consequently:

\begin{itemize}
  \item \textbf{``$p = 1.228$ supports H1 over H2'' is withdrawn.} One local
  effective order on one branch cannot distinguish asymptotic first order from
  pre-asymptotic behaviour of a second-order scheme. The H1/H2 question is
  open.
  \item \textbf{The $1$~N error decomposition ($1.21$~kg basin, $0.35$~kg
  mesh) is withdrawn.} It subtracted a Richardson limit extrapolated on the
  $171$-switch branch from a production result on the $177$-switch branch.
  Both reviewers called this incoherent, and they are right: it is arithmetic,
  not an identifiable decomposition.
  \item \textbf{The dominance ratios ($\sim\!100\times$, $\sim\!4\times$)
  are withdrawn.} The $10$~N figure rested on a ``within-branch mesh error'' of
  $0.009$~kg whose two levels have $19$ and $20$ switches---not within-branch
  at all. The $1$~N figure rested on the withdrawn decomposition.
  (A later attempt to re-establish a dominance ratio of $4.4$ at $10$~N was
  itself withdrawn---see Section~\ref{sec:corrections2}. No supported dominance
  ratio exists.)
  \item \textbf{``The production meshes are broadly adequate'' is
  withdrawn.} No branch-matched convergence estimate exists for any row.
\end{itemize}

\paragraph{What survives, and why.} The basin results are unaffected, because
they never depended on any convergence argument. They compare solutions of the
\emph{same} nonlinear program on the \emph{same} grid, all machine-feasible,
and observe that some have higher final mass than the certified one. That is a
direct observation, not an extrapolation.

\paragraph{The process lesson.} The Richardson restriction had been written
down, reviewed, and agreed---and was then violated by the author of the rule,
because it lived in a header comment while the arithmetic did not enforce it.
\texttt{mesh\_order} now takes an explicit per-difference validity mask and
refuses any window spanning a branch change; \texttt{mesh\_collect} builds
that mask from the switch counts. A rule that is not executable is a rule that
will be broken.

\section{Introduction}
\label{sec:intro}

\subsection{What is already verified, and what is not}

Our low-thrust transfer campaigns solve minimum-fuel optimal control problems
by direct collocation: the continuous problem is discretized on a mesh of
nodes, the dynamics are enforced as algebraic \emph{defect} constraints between
adjacent nodes, and the resulting large sparse nonlinear program (NLP) is
handed to an interior-point solver. Three layers of verification are already in
place, and it is worth being precise about what each one does and does not
establish.

\begin{enumerate}
  \item \textbf{The solver's own convergence test.} IPOPT reports
  \texttt{Solve\_Succeeded} when its primal and dual infeasibilities fall below
  tolerance. This establishes that a stationary point of the NLP was found.
  \item \textbf{Feasibility of the discretization.} We check that the maximum
  defect is at machine precision (typically $10^{-14}$ or smaller). This
  establishes that the discrete trajectory satisfies the discrete dynamics.
  \item \textbf{First-order optimality.} Our \texttt{foc\_check} layer assembles
  the Lagrangian gradient by automatic differentiation and verifies KKT
  stationarity, the Pontryagin primer-vector alignment, transversality, and the
  sign law on the multipliers.
\end{enumerate}

Every one of these is a statement about the \emph{nonlinear program}. None is a
statement about the continuous optimal control problem. A solution can pass all
three and still differ from the true continuous optimum, because the mesh is
too coarse to represent it. Nothing in the repository measured that gap. This
document records the study that does.

\subsection{The two questions, which turn out to be different}

The study began with one question and ended with two, and keeping them separate
is the single most important organizational point in this document.

\begin{description}
  \item[Q1 (discretization).] Is the production mesh fine enough? Equivalently:
  if we refined it indefinitely, how much would the answer move?
  \item[Q2 (basin selection).] Is the solution we certified even the best
  solution \emph{of the mesh we already have}?
\end{description}

These are independent. A row can be perfectly mesh-converged and still sit in
an inferior local minimum, because the minimum-fuel bang-bang problem is
non-convex. Conversely a row can be in the best available basin and still be
under-resolved. The study was designed to answer Q1; in the course of answering
it we were forced to answer Q2, and Q2 turned out to carry the larger error.

\section{The mesh, and what mesh convergence means}
\label{sec:mesh}

\subsection{Our meshes are not uniform in time}

A point that recurs throughout this document: none of our meshes is uniform in
physical time, and this is deliberate.

The Earth-centred campaign uses modified equinoctial elements with the true
longitude $L$ as the independent variable, so a uniform grid in the normalized
variable $\sigma \in [0,1]$ is uniform in longitude, not in time. The
cislunar campaigns use a Sundman regularization, replacing time $t$ by a
variable $\tau$ with
\begin{equation}
  \frac{dt}{d\tau} = \kappa(r) = r_1^{\,p}, \qquad p = 1.5,
  \label{eq:sundman}
\end{equation}
where $r_1$ is the distance to the primary. Physical time is carried as an
additional state. A uniform grid in $\tau$ therefore clusters nodes near
perigee automatically, which is exactly where the dynamics are stiff.

The practical consequence is that the local step size in physical time varies
enormously across a single transfer. Measured on the Earth $10$~N row, the
physical step varies by a factor of $36$ from its smallest to its largest value
even though $\sigma$ is perfectly uniform. Any statement of the form ``the mesh
spacing is $h = 1/N$'' is therefore wrong for our problems, and any tolerance
built from a \emph{global} step size will be simultaneously too tight in one
region and too loose in another. Throughout this study we report physical local
step sizes and use per-node tolerances.

\subsection{The refinement experiment}

The logic of a mesh-convergence study is simple. Take a converged solution on
the production mesh. Re-solve the identical continuous problem---same boundary
conditions, same transfer time, same thrust---on meshes with $2\times$,
$4\times$ and $8\times$ as many nodes. Then look at how a quantity of interest,
say the final mass $\mf$, behaves across the ladder.

If $\mf$ barely moves, the production mesh was already fine enough. If it moves
substantially at every refinement, we are not converged and the production
number carries an error we have not accounted for.

Two refinements of that basic logic matter, and both were learned the hard way.

\paragraph{It is the \emph{rate}, not the size, of the change that matters.}
Knowing that the answer moved by $\delta$ between two meshes tells you very
little on its own. What tells you something is how $\delta$ shrinks as the mesh
refines, because that rate is what lets you extrapolate to the answer you would
obtain at infinite resolution. This is the subject of
Section~\ref{sec:richardson}. Two mesh levels give you a change; three or four
give you a rate and hence an error bar. This is why the study specifies four
levels rather than two.

\paragraph{The comparison is only valid on a single solution branch.} The entire
argument silently assumes that the thing being compared across meshes is the
\emph{same solution}, merely represented more finely. If the solver lands in a
different local minimum at one of the levels, then the differences being
tabulated are differences between unrelated answers and the extracted ``rate''
is meaningless. As Section~\ref{sec:results} records, this is not a hypothetical
concern: it happened in two of our three rows.

\subsection{Warm chaining, and the risk it carries}

Each level of the ladder is solved starting from the previous, coarser
solution, interpolated onto the finer grid. This is deliberate. Our basins are
separated by very small perturbations---an earlier campaign study documented
the $10$~N case flipping between $19$ and $24$ switches on a relative change of
$2\times10^{-5}$ in the transfer time---so cold-solving each level
independently would risk landing on unrelated branches and comparing apples to
oranges.

The cost of chaining is the mirror image of its benefit: a level can inherit an
unconverged or inferior structure from its predecessor and still pass every
gate, because passing a gate proves only that the level is stationary for
\emph{its own} NLP, not that all levels lie on the same continuous branch.
Detecting that requires an independent check, described in
Section~\ref{sec:basin}.

\section{Richardson extrapolation and the observed order}
\label{sec:richardson}

This section explains the method referred to elsewhere as ``Richardson,''
because it is the machinery that converts a sequence of mesh solutions into an
error bar.

\subsection{The idea}

Suppose a computed quantity $f_h$, obtained on a mesh of characteristic
spacing $h$, approximates an exact value $f_0$ with an error that behaves like
a power of $h$:
\begin{equation}
  f_h = f_0 + C h^p + \mathcal{O}(h^{p+1}),
  \label{eq:expansion}
\end{equation}
for some constant $C$ and some exponent $p > 0$ called the \emph{order of
accuracy}. Equation~\eqref{eq:expansion} is an assumption, not a theorem, and
Section~\ref{sec:richardson-limits} discusses when it fails for us.

Evaluate the quantity on three meshes whose spacings shrink by a constant
refinement ratio $r$ (we use $r=2$, doubling the node count each time):
$h$, $h/r$, $h/r^2$, giving values $f_1, f_2, f_3$. Substituting
\eqref{eq:expansion} and taking differences eliminates $f_0$:
\begin{align}
  f_1 - f_2 &= C h^p \left( 1 - r^{-p} \right), \\
  f_2 - f_3 &= C h^p r^{-p} \left( 1 - r^{-p} \right),
\end{align}
so the ratio of successive differences depends only on $r$ and $p$:
\begin{equation}
  \frac{f_1 - f_2}{f_2 - f_3} = r^{p}
  \qquad\Longrightarrow\qquad
  \boxed{\;p = \frac{\ln\!\left| (f_1-f_2)/(f_2-f_3) \right|}{\ln r}\;}
  \label{eq:order}
\end{equation}
This $p$ is the \emph{observed order of accuracy}. It is measured, not assumed,
and comparing it against the order the scheme is supposed to achieve is itself
a diagnostic.

\subsection{Extrapolating to the limit}

With $p$ in hand, \eqref{eq:expansion} can be solved for $f_0$:
\begin{equation}
  f_0 \;\approx\; f_3 + \frac{f_3 - f_2}{r^{p} - 1}.
  \label{eq:richardson}
\end{equation}
This is the \emph{Richardson extrapolate}: an estimate of the answer at
infinite resolution, obtained from finite-resolution runs. Its practical value
here is that $|f_0 - f_{\text{production}}|$ is an \emph{error bar on the
production number}---precisely the quantity we lacked.

A worked instance from Section~\ref{sec:results}: the $1$~N row gives
$\mf = 1371.587,\ 1371.913,\ 1372.052$~kg on the $2\times$, $4\times$ and
$8\times$ meshes. The differences are $+0.326$ and $+0.139$, whose ratio is
$2.34$, so $p = \ln(2.34)/\ln 2 = 1.23$. Extrapolating,
$f_0 \approx 1372.052 + 0.139/(2^{1.23}-1) = 1372.156$~kg. The finest mesh is
thus about $0.10$~kg short of the limit.

\subsection{Where it must not be used}
\label{sec:richardson-limits}

Equation~\eqref{eq:expansion} assumes a smooth asymptotic error expansion in
$h$. Several quantities in our problems violate that assumption, and applying
the formula to them yields confident nonsense. Our implementation therefore
maintains an explicit permitted/forbidden classification and \emph{refuses} to
report an order where it is forbidden.

\begin{description}
  \item[Forbidden: integer counts.] The number of thrust switches is an
  integer. It has no asymptotic expansion; it must \emph{stabilize}, and if it
  does not, that is a finding about the solution's topology rather than a
  convergence rate.
  \item[Forbidden: solver diagnostics.] KKT residuals, multiplier sign
  percentages and similar readouts measure the NLP and the solution process,
  not the continuous problem. They are tracked and reported but never given a
  convergence verdict alongside physical quantities.
  \item[Forbidden: anything measured across a branch change.] If the topology
  changes between levels, the differences are not samples of one smooth
  expansion.
  \item[Forbidden: quantities pinned by constraints.] The terminal element
  error, for instance, is enforced as an equality in the NLP, so it reads
  solver tolerance by construction ($\sim 10^{-35}$ in our runs) and can never
  measure discretization error. Classifying it as a physical observable would
  earn it a meaningless ``converged'' verdict.
\end{description}

Two further guards are built into the estimator. If successive differences flip
sign or stop shrinking, the sequence is not settling and no order is reported.
And with four levels rather than three, two overlapping three-level windows can
be computed; their mutual agreement is the evidence that an asymptotic regime
has actually been reached, which no single window can supply.

\section{Basin selection}
\label{sec:basin}

\subsection{Why a non-convex problem has this failure mode}

The minimum-fuel problem is non-convex. At the end of the
energy-to-fuel homotopy the throttle is bang-bang: it takes only the values $0$
and $1$, switching between them at a finite set of times. Solutions are
therefore organized into families distinguished by discrete structure---the
number of thrust arcs and their approximate placement. Within a family the
solution varies smoothly; between families it does not. An interior-point
solver converges to a stationary point in whichever family its initial guess
lies nearest, and it has no mechanism for discovering that a different family
would have been better.

The consequence is that ``converged'' and ``optimal'' are different claims.
Our certification establishes the first. Nothing in the pipeline established
the second.

\subsection{Distinguishing a basin effect from a resolution effect}

These two effects are easy to confuse, because both appear as ``the answer
changed when I refined the mesh.'' The test that separates them is a
\emph{down-projection}:

\begin{quote}
Take the solution obtained on the fine mesh. Interpolate it \emph{down} onto
the original production mesh. Re-solve there. If the fine-mesh structure
survives at production resolution and is better, then the fine mesh did not
resolve anything the coarse mesh could not represent---it merely escaped to a
different basin, and that basin was reachable at the original resolution all
along.
\end{quote}

This is a clean discriminator because the comparison is made on an identical
grid, so nothing about it can be a discretization effect.

\subsection{The basin probe}

Generalizing the down-projection gives the tool used in
Section~\ref{sec:results}: re-solve a certified row on \emph{its own} mesh from
a set of different seeds, and compare. Because the grid never changes, any
improvement is attributable to basin selection alone. The seeds used are:

\begin{itemize}
  \item the certified solution itself, as a baseline;
  \item down-projections of each finer mesh-ladder level;
  \item the throttle profile shifted by a small number of nodes---a structured
  nudge that moves the \emph{switch times}, which is the coordinate the basins
  actually differ in, and which proved far more effective than random
  perturbation (at $\varepsilon = 0$ a random perturbation typically falls
  straight back into the same basin);
  \item for the cislunar campaigns, a re-solve with a looser barrier
  initialization, which is not a seed change at all but a different solver
  path.
\end{itemize}

Throughout, \emph{better} means higher final mass: minimum fuel maximizes
$\mf$, so a higher mass at machine-tight feasibility is strictly an
improvement, never a failure.

\section{Results}
\label{sec:results}

\subsection{What was run}

The Earth-centred ladder (``Tier A'') re-solved three certified rows---thrust
levels $10$~N, $2.5$~N and $1$~N---at node-count multipliers
$1, 2, 4, 8$, warm-chained coarse to fine. The $1$~N row at $8\times$ carries
$13{,}920$ intervals. All twelve levels converged, with defects between
$5.5\times10^{-15}$ and $9.7\times10^{-14}$; total wall time was $80.3$
minutes. All artifacts are written under a \texttt{MESH\_} prefix into an
isolated directory so that a study solve cannot overwrite a production cache.

\subsection{Tier A: the mesh ladder}

\begin{table}[h]
\centering
\small
\begin{tabular}{lrrrrl}
\toprule
Row & $N$ ($1\times\!\to\!8\times$) & $\mf$ per level (kg) & differences & $p$ windows & verdict \\
\midrule
$10$~N  & $193\to1544$   & $1377.101\,/\,1377.110\,/\,1377.118\,/\,1377.875$
        & $+0.0086,\ +0.0077,\ +0.757$ & $0.16,\ -6.62$ & not converged \\
$2.5$~N & $708\to5664$   & $1370.157\,/\,1371.612\,/\,1372.610\,/\,1375.794$
        & $+1.46,\ +1.00,\ +3.18$      & $0.54,\ -1.67$ & not converged \\
$1$~N   & $1740\to13920$ & $1370.603\,/\,1371.587\,/\,1371.913\,/\,1372.052$
        & $+0.984,\ +0.326,\ +0.139$   & $1.60,\ 1.23$  & \textbf{converged} \\
\bottomrule
\end{tabular}
\caption{Tier A mesh ladder. Switch counts: $10$~N $19/20/20/18$;
$2.5$~N $76/75/75/75$; $1$~N $173/171/171/171$.}
\end{table}

\paragraph{Order estimates, after enforcing branch consistency.} Applying the
rule of Section~\ref{sec:richardson-limits} in code---no window may span a
change in switch count---leaves very little standing:

\begin{center}
\small
\begin{tabular}{lccc}
\toprule
Row & switch counts & valid windows & $p$ \\
\midrule
$10$~N  & $19/20/20/18$    & $0$ & none \\
$2.5$~N & $76/75/75/75$    & $1$ & $-1.67$ (differences grew) \\
$1$~N   & $173/171/171/171$& $1$ & $1.228$ \\
\bottomrule
\end{tabular}
\end{center}

\textbf{No row has an established asymptotic order.} The $1$~N value of
$1.228$ is a single local effective order on the $171$-switch branch. It is
below $2$, but one window cannot distinguish asymptotic first order from
pre-asymptotic behaviour of a second-order scheme, and it cannot support a
statement about a ``trend.'' Its accompanying Richardson value, $1372.156$~kg,
is a limit for that branch only.

The H1/H2 question therefore remains open, which is worth stating plainly
because the first draft claimed to have settled it. Before the study, two
external reviewers disagreed about what to expect. One argued for
$p \approx 1$: on a mesh whose breakpoints
are not aligned with the switches, one switch-crossing interval carries the
wrong throttle over a time of order $h$ against an order-one jump in the
right-hand side, and that single interval's $\mathcal{O}(h)$ error dominates
every smooth arc's $\mathcal{O}(h^2)$. The other argued for $p \approx 2$,
on the grounds that the throttle is a continuous NLP variable so an
intermediate node value can encode a sub-interval switch location, and that
errors of alternating sign across many switches partially cancel in an
integrated quantity. We recorded both as competing falsifiable predictions before running, precisely
so that the result would discriminate rather than confirm whichever had been
written down first. The measurement does not discriminate: one window on one
branch is not enough. Settling it requires continuing a \emph{single} branch
to $16\times$ and beyond, and verifying branch identity by matched switch
times rather than by switch count alone.

\paragraph{Two rows did not converge.} Both $10$~N and $2.5$~N had their
$8\times$ difference \emph{grow} rather than shrink. At $10$~N the switch count
also dropped from $20$ to $18$ and the level took $179.5$~s against $11.7$~s at
$4\times$---fifteen times longer, far above the empirical $N^{1.88}$ scaling,
because the solver was doing something other than refining. At $2.5$~N the
switch count did \emph{not} change, yet the mass still jumped by $3.18$~kg.

\subsection{The down-projection test}

Applying the discriminator of Section~\ref{sec:basin} to the $10$~N row settles
what happened.

\begin{table}[h]
\centering
\small
\begin{tabular}{lrrr}
\toprule
On the same $193$-node production mesh & $\mf$ (kg) & switches & $\dV$ (km/s) \\
\midrule
seeded from the row itself (certified) & $1377.101235$ & $19$ & $1.676631$ \\
seeded from the $8\times$ solution     & $\mathbf{1378.116718}$ & $\mathbf{18}$ & $\mathbf{1.662173}$ \\
\bottomrule
\end{tabular}
\caption{The $18$-switch structure survives projection onto the production
mesh and is $1.015$~kg better there. Both solves reported
\texttt{Solve\_Succeeded} with defects below $10^{-14}$.}
\end{table}

The $8\times$ level did not resolve something the coarse mesh could not
represent. It escaped a basin, and the better basin is reachable at production
resolution. Asking whether the change was ``real or a solver artifact'' poses a
false choice: the $8\times$ solution is a genuine converged optimum, and it is
not a mesh phenomenon.

\subsection{The basin probe on the Earth rows}

\begin{table}[h]
\centering
\small
\begin{tabular}{llrrrr}
\toprule
Row & best seed & $\mf$ certified (kg) & $\mf$ best (kg) & gain (kg) & switches \\
\midrule
$10$~N  & down-proj.\ $8\times$ & $1377.101$ & $1378.117$ & $+1.015$ & $19 \to 18$ \\
$2.5$~N & down-proj.\ $8\times$ & $1369.791$ & $1375.762$ & $+5.971$ & $76 \to 69$ \\
$1$~N   & down-proj.\ $4\times$ & $1370.357$ & $1371.811$ & $+1.454$ & $171 \to 177$ \\
\bottomrule
\end{tabular}
\caption{Every certified Earth row is beaten on its own mesh, and in every case
the winner also has lower $\dV$.}
\end{table}

Three observations sharpen this.

\begin{enumerate}
  \item \textbf{A trivial perturbation suffices.} At $10$~N, shifting the
  throttle profile by a \emph{single node} reaches an $18$-switch solution
  $0.224$~kg better than certified. The better basin is not exotic; it was
  simply never sought.
  \item \textbf{Switch count is not a reliable branch indicator.} The $2.5$~N
  ladder kept $75$ switches throughout, which made its jump look like it might
  not be a topology change. Down-projected onto the production mesh it lands at
  $69$ switches and $+5.97$~kg. It was a different branch all along.
  \item \textbf{Re-solving a row from its own seed does not reproduce its
  cached value.} For $2.5$~N the cache reads $1369.791$~kg, the probe's
  baseline gives $1370.296$, and the ladder's $1\times$ level gives $1370.157$.
  These differ only in trivial seed handling---whether the seed's thrust
  direction is re-normalized to unit length and its throttle clipped into
  $[0,1]$. A seed change at the $10^{-12}$ level moves the answer by half a
  kilogram. ``The certified value'' is therefore more precisely ``the value
  this particular solve path happened to reach.''
\end{enumerate}

\subsection{Reconciling the two questions}

A first draft tabulated ``basin gain'' against ``residual mesh error'' and
concluded that basin selection dominates by factors of four to one hundred.
That table has been withdrawn (Section~\ref{sec:corrections}): it compared a
discrete jump between distinct local minima against a continuum residual, and
worse, both of its mesh-error entries were invalid---the $10$~N figure spanned
a $19\to20$ switch change, and the $1$~N figure came from the withdrawn
cross-branch subtraction.

What can be said without any convergence argument at all is this:

\begin{quote}
For all three Earth production problems, changing only the initial guess found
feasible solutions better by $1.015$, $5.971$ and $1.454$~kg respectively, each
with lower $\dV$, on the unmodified production mesh.
\end{quote}

This is an observation about the nonlinear programs we actually solve, not an
extrapolation, and it is unaffected by the retractions. It establishes that
basin search is an immediately higher-yield intervention than mesh refinement.
It does \emph{not} establish a quantitative error budget in which basin error
is shown to exceed discretization error, because no valid branch-matched
discretization error has been measured for any row.

Obtaining one is straightforward in principle and is the natural next
experiment: take the \emph{best known} production solution for a row, continue
\emph{it} up the ladder holding its branch, verify branch identity by matched
switch times rather than by count, and extrapolate along that single branch.

\subsection{Cislunar campaigns: tulip and ELFO}
\label{sec:cr3bp}

\paragraph{Tulip flagship: a null result, with a caveat that matters.} All six
seeds---the certified solution, throttle shifts of $-8$, $-1$, $+1$ and $+8$
nodes, and a re-solve under a looser barrier initialization---returned to
$\mf = 0.84906583$ with $25$ switches, agreeing to $10^{-10}$ or better. The
flagship survives this probe, in sharp contrast to every Earth row, where a
single-node shift was sometimes enough to find a better basin.

The caveat is that \emph{this seed family is the weak one}. No mesh ladder has
been run for the cislunar campaigns, so the down-projection seeds are
unavailable---and those were the winning seed at every single Earth row. Both
external reviewers made the same point independently: a null result under
perturbations that are known to be weaker than the ones that worked elsewhere
establishes local robustness, not optimality. Concluding that this row is
robustly optimal requires subjecting it to the same adversary: run a cislunar
mesh ladder, down-project each level, and verify the flagship resists.

\paragraph{ELFO: not yet run.} The first attempt failed on all seeds because
the probe called the wrong solver. The ELFO campaign does not use the tulip's
fixed-$t_f$, single-primary, $8$-state formulation; it uses a free-$t_f$
two-primary solver with a $9$-state vector (the ninth being a time-scaling
slack) and a two-primary Sundman clock
$\kappa = (r_1^{-q} + (r_2/D)^{-q})^{-p/q}$, needed because a single-primary
Earth clock cannot resolve the lunar capture leg. The probe assumed one CR3BP
solver where there are two.

\section{Consequences and open items}
\label{sec:consequences}

\paragraph{Published numbers are pessimistic.} The Earth rows understate
achievable final mass by $0.07\%$ ($10$~N), $0.11\%$ ($1$~N) and $0.44\%$
($2.5$~N), each with lower $\dV$ as well, so these are not trade-offs. One
qualification: at $2.5$~N, re-solving the row from its own solution---differing
only in whether the seed's thrust direction was renormalized and its throttle
clipped---moves the answer by $0.5$~kg. A gain must be read against that
per-row reproducibility scale, not against a single global one. The $2.5$~N
gain of $5.971$~kg clears its own scale by an order of magnitude; the $10$~N
gain of $1.015$~kg sits on a row whose baseline reproduced exactly, so its
local scale is far smaller. The document previously described the perturbation
as ``a $10^{-12}$ seed change''; that was not measured, and projection onto the
unit-direction and box constraints can amplify a small raw change, so the
characterization is withdrawn.

\paragraph{A basin-search step belongs in the certification path.} At present
``certified'' means converged, feasible and first-order optimal---not best
found. The probe costs three to fifteen minutes per row.

\paragraph{Switch counts should not be quoted as properties of the problem.}
They vary with basin and with mesh: $18/19/20$ at $10$~N, $171/173/177$ at
$1$~N. They are properties of a particular solve.

\paragraph{Open items, in priority order.}
\begin{enumerate}
  \item \textbf{A branch-matched ladder.} Continue the \emph{best known}
  production solution for one row up the ladder holding its branch, verifying
  branch identity by matched switch times rather than switch count. Without
  this there is no valid mesh-error estimate for any row, and the study's
  first question remains unanswered.
  \item \textbf{Extend that ladder to $16\times$} to settle H1 versus H2. One
  window cannot.
  \item \textbf{Port the probe to the ELFO solver} and run it.
  \item \textbf{A cislunar mesh ladder}, so the tulip null can be tested
  against the down-projection seeds that beat every Earth row.
  \item \textbf{A per-row reproducibility scale}, measured rather than
  assumed, so that a reported gain can be read against it.
  \item Test the falsifiable claim that switch-aligned refinement recovers
  second order---only if it makes switches true retained breakpoints.
\end{enumerate}


\section{Phase 3: continuous Pontryagin residuals}
\label{sec:phase3}

\subsection{The redirection, and why it works}

Every difficulty in Sections~\ref{sec:mesh}--\ref{sec:results} has one root.
Mesh convergence is a \emph{comparative} measurement: it needs two or more
solutions, and therefore a proof that they are the same solution merely
represented more finely. That proof is what could not be supplied---each ladder
crossed a change in the discrete switch structure, and external review found
that no tightening of criteria would fix it, because the underlying
continuation cannot hold a branch.

The redirection asks an \emph{absolute} question of a \emph{single} solution:

\begin{quote}
Take the converged trajectory, read it as a continuous object, and ask whether
it satisfies the continuous necessary conditions of the indirect (Pontryagin)
formulation---and to what residual.
\end{quote}

Nothing is compared, so no branch-identity question arises. It also closes a
gap external review called blocking: \texttt{Solve\_Succeeded} plus a nodal
defect bound verifies only that the solver solved \emph{the transcription}, and
nothing in the repository asked whether the result approximates a continuous
extremal.

\subsection{What we hoped to learn}

Three things, in decreasing order of confidence that they were obtainable:

\begin{enumerate}
  \item \textbf{Whether the discrete solution is close to a continuous
  extremal at all}, and by how much it misses. This is the question the defect
  was silently assumed to answer and does not.
  \item \textbf{Where the error lives.} A residual computed per interval is a
  \emph{map}, not a scalar, so it localizes the difficulty---which is far more
  actionable than a single number.
  \item \textbf{Whether the bang-bang structure is the dominant error source},
  as hypothesis H1 asserts. This had been argued in both directions by external
  reviewers and never measured.
\end{enumerate}

\subsection{Method, and the two constraints that shape it}

The conditions checked are the primal dynamics, the costate dynamics, the
boundary values, transversality, the minimum condition $u = \arg\min_u H$,
constancy of the conserved quantity, switching regularity, and the multiplier
sign conditions. Of these the costate ODE was checked by \emph{nothing} before,
and the minimum condition only at nodes.

\paragraph{Propagation is per interval, never global.} The costate dynamics are
exponentially unstable; a single forward shot from $\lambda(0)$ across a
many-revolution transfer diverges and measures nothing. Each interval is
instead integrated from the solution's own values at its left node and compared
at its right node, so nothing accumulates.

\paragraph{The dynamics are the transcription's own.} The right-hand side is
reused directly and $\partial H/\partial x$ taken from it by automatic
differentiation. The check must be independent in \emph{integration}, not in
the definition of the physics: a re-derived right-hand side would test the
author's algebra rather than the mesh. A self-validating gate recomputes the
transcription's own defects from the reused right-hand side and aborts if they
do not reproduce; a residual computed with the wrong dynamics is worse than no
residual, because it looks like evidence.

\subsection{P3.1: validation against an analytic extremal}

The double integrator $\min\int_0^2 u\,dt$ with $\dot x = v$, $\dot v = u$,
$u\in[0,1]$, $x(2)=1$, $v(2)$ free has the closed-form extremal
$t_s = 2-\sqrt2$, $\lambda_x = -1/\sqrt2$, $H = 1-\sqrt2$.

\begin{table}[h]
\centering\small
\begin{tabular}{lr}
\toprule
check & result \\
\midrule
self-gate, both arcs & $1.7\times10^{-17}$ / $8.3\times10^{-17}$ \\
primal residual $R_x$ & $\max\ 1.5\times10^{-16}$ \\
costate residual $R_\lambda$ & $\max\ 1.1\times10^{-15}$ \\
Hamiltonian & $-0.414213562373$, analytic to 12 digits, CoV $1.5\times10^{-15}$ \\
unaligned grid, straddling interval & $\mathbf{1.71\times10^{-3}}$ \\
unaligned grid, elsewhere & $3.3\times10^{-16}$ \\
\bottomrule
\end{tabular}
\end{table}

Both halves are required. The first shows the checker is \emph{correct}; the
second that it is \emph{sensitive}, with thirteen orders of separation and the
peak at the right index. A checker passing the first but missing the second
would be blind to the only error the study cares about.

A structural lesson emerged here rather than being assumed: at a bang-bang
switch the control is genuinely double-valued, so a single shared node cannot
carry both $u=1$ and $u=0$. A \emph{true} retained breakpoint is a
\textbf{duplicated} node---which is exactly what a multi-phase formulation
provides, and is the concrete mechanism behind the switch-alignment idea.

\subsection{P3.2--P3.3: the three earth rows}

All self-gates reproduce the solver's own reported defect
($6.245$ vs $6.273\times10^{-15}$ at $10$~N; $8.266\times10^{-14}$ exactly at
$1$~N). The conserved quantity for this transcription is the \emph{time
costate}---$H_\sigma$ is not constant, since the Gauss equations carry $\sigma$
explicitly through $L = \pi + \sigma\Delta L$---and $\lambda_t$ holds to
CoV $2.3\times10^{-8}$ in every row.

\paragraph{The headline gap.} The median discrete defect is
$5.0\times10^{-16}$; the median continuous residual is $3.7\times10^{-4}$. The
solution satisfies the equations it was \emph{solved against} roughly
$7\times10^{11}$ times better than the continuous equations we care about.

\paragraph{Finding 3.1 (a correction).} The raw $R_\lambda$ reaching $13.9$ was
flagged as possibly alarming. It is an \emph{absolute} norm and costates are
sensitivities with $|\lambda|$ of $24$--$264$. Normalized, the median relative
costate residual is $4.6\times10^{-6}$, $7.7\times10^{-6}$ and
$1.4\times10^{-6}$ across the three rows; one interval anywhere reaches 22\%.

\paragraph{Finding 3.2: the peak is at apogee.} Median $R_x$ in the
apogee-most decile against the perigee-most decile is $16.5\times$ ($10$~N),
$19.2\times$ ($2.5$~N) and $18.2\times$ ($1$~N). The mechanism: the mesh is
uniform in true \emph{longitude}, and near apogee the angular rate is slow, so
a fixed $\Delta L$ spans a long $\Delta t$. The longest time-steps sit at
apogee, and that is where a trapezoid hurts most.

\paragraph{Finding 3.3: it tracks neither switches nor revolutions.}

\begin{table}[h]
\centering\small
\begin{tabular}{lrrrrrr}
\toprule
row & revs & switches & sw/rev & nodes/rev & median $R_x$ & \textbf{median $R_x/h^3$} \\
\midrule
$10$~N  & 7.33  & 19  & 2.59 & 26.3 & $3.73\times10^{-4}$ & $\mathbf{0.171}$ \\
$2.5$~N & 27.82 & 76  & 2.73 & 25.4 & $4.61\times10^{-4}$ & $\mathbf{0.171}$ \\
$1$~N   & 69.18 & 171 & 2.47 & 25.2 & $4.65\times10^{-4}$ & $\mathbf{0.163}$ \\
\bottomrule
\end{tabular}
\caption{The median residual is flat while revolutions, switches and node count
each vary by a factor of nine. Normalized by $h^3$ it is constant to 5\%.}
\end{table}

That is precisely the local truncation error of the trapezoidal rule,
$\mathcal{O}(h^3)$ per step, with a coefficient set by the \emph{orbital
dynamics} rather than by thrust level or control structure. The campaign fixes
nodes per revolution, so the same steps recur at the same orbital phases in
every row and the same residual follows.

\subsection{What Phase 3 established}

\begin{itemize}
  \item The pre-registered prediction that residuals would concentrate at
  switch intervals \textbf{failed}: switch cells are only $3.6\times$ the
  interior median at $10$~N, and the largest residual is an interior cell.
  Findings 3.2 and 3.3 say what happens instead.
  \item H1's mechanism---one switch-crossing interval swamping the smooth
  arcs---is not what drives the error in these problems.
  \item Judged by \emph{error equidistribution} rather than by angular
  resolution, the mesh over-resolves perigee and under-resolves apogee. That is
  a concrete target for node placement, and it points \emph{away} from
  switch-aligned refinement as the lever.
\end{itemize}

\subsection{What Phase 3 does not deliver}

\paragraph{A mass error bar.} The residual is local and per interval. It
establishes that the mesh is mis-allocated; it does not say what that costs in
kilograms. The conversion is the first item in Section~\ref{sec:next}.

\paragraph{Anything about basins.} The Pontryagin conditions are
\emph{necessary}: every local extremal satisfies them. The $10$~N solution that
is $1.015$~kg worse than its neighbour passes this check cleanly.

\subsection{Step 1: converting the residual to kilograms}

The residual is local and dimensionless; the question was always in kilograms.
The bridge is the multiplier. The discrete solution satisfies its own defect
constraints exactly, but measured against the exact flow each interval carries
a signed residual $d_k$, so the discrete solution is the exact solution of a
problem whose defect constraints are displaced by $d_k$. The multiplier of
defect $k$ is by definition $\partial J^*/\partial(\text{defect }k)$, hence to
first order
\begin{equation}
  \Delta J \;\approx\; \sum_k \lambda_k^{\!\top} d_k .
  \label{eq:sens}
\end{equation}
No fresh adjoint propagation is needed---the nonlinear program already computed
these sensitivities, which also avoids the backward instability that forced
per-interval propagation in the first place.

\begin{table}[h]
\centering\small
\begin{tabular}{lrrrrr}
\toprule
row & $\mf$ (kg) & implied $|\Delta m|$ & no-cancel bound & cancellation & relative \\
\midrule
$10$~N  & 1377.101 & 1.237 & 3.517 & 2.8 & $9.0\times10^{-4}$ \\
$2.5$~N & 1370.296 & 5.714 & 8.264 & 1.4 & $4.2\times10^{-3}$ \\
$1$~N   & 1370.798 & 4.046 & 6.634 & 1.6 & $3.0\times10^{-3}$ \\
\bottomrule
\end{tabular}
\end{table}

Reporting the signed sum beside the sum of magnitudes measures how much the
per-interval errors cancel, and that ratio is a direct test of H2, which held
that alternating-sign errors across many switches would partially cancel in an
integrated quantity. Measured, cancellation removes only a factor of 1.4--2.8,
and it is \emph{weakest} on the rows with the \emph{most} switches---the
opposite of the argument's prediction. H2's cancellation mechanism does not
operate here.

\subsection{Step 2: mesh redistribution, and the closed loop}

Equation~\eqref{eq:sens} is a first-order estimate; whether it is quantitatively
right is a separate question, and the way to answer it is to predict a change
and then make it. Given the measured $E \propto h^3$, the redistribution that
minimizes $\sum h^3$ at fixed node count and fixed horizon is $h$ constant---a
mesh uniform in physical \emph{time} rather than in longitude. Predictions were
computed and recorded \emph{before} the re-solve.

\begin{table}[h]
\centering\small
\begin{tabular}{clll}
\toprule
\# & prediction & measured & \\
\midrule
1 & $\sum h^3$ drops $2.68\times$ & $2.66\times$ & \textbf{hit} \\
2 & residual falls by the same factor & median $1.47\times$, max $8\times$ \emph{worse} & \textbf{miss} \\
4 & $\Delta m \approx 0.776$~kg & $+0.146$~kg & \textbf{miss}, $5.3\times$ \\
5 & switch count stays 19 & $19 \to 19$ & \textbf{hit} \\
\bottomrule
\end{tabular}
\end{table}

The step-size range collapsed from $36.2\times$ to $1.6\times$, and the switch
count held---so this is a \emph{single-branch} comparison, and the objection
that defeated Phase~2 does not apply to it.

\paragraph{The sign was inverted, and is now settled.} Step~1's estimate implied
the discrete solution \emph{overstates} achievable mass, so a better mesh should
lower $\mf$. It \emph{rose}, $1377.101235 \to 1377.246863$. The discrete
solution understates achievable mass and improving the mesh raises it, agreeing
with the direction seen in the ladder. This was settled by perturbation rather
than by arguing about multiplier conventions.

\paragraph{The first-order estimate overpredicts by a factor of five.} Predicted
$0.776$~kg against a measured $0.146$~kg. Equation~\eqref{eq:sens} is an
order-of-magnitude indicator, not a calibrated error bar, and should not be
leaned on as one.

\subsection{The attempted error bar, and why it does not stand}

Two solutions on the same node count, differing only in node placement, with
$\sum h^3$ of $2.668$ and $1.003$ and final masses $0.1456$~kg apart. Taking
$E = k\sum h^3$ gives $k = 0.0874$~kg per unit and a production-mesh error of
$0.233$~kg. \textbf{That number was reported as the study's deliverable and is
withdrawn}; Section~\ref{sec:corrections2} gives the reasoning.

What survives the withdrawal is weaker but sound. Both solutions understate the
achievable mass---this is the direction established by the sign resolution
above---so their difference bounds the worse one from below:

\begin{quote}
The $10$~N production mesh carries \textbf{at least} $0.1456$~kg of mass
discretization error, conditional on the two solutions lying on one branch.
\end{quote}

That is a lower bound, not a point estimate, and it still carries a conditional.

\section{Corrections after the Phase 3 review}
\label{sec:corrections2}

The Phase 3 results were sent for external review with the measurements and
figures attached. Both reviewers returned \textbf{not defensible}, naming the
same weakest link.

\paragraph{The model was falsified by our own data.} $E = k\sum h^3$ factors a
single constant out of $\sum_i C_i h_i^3$, which assumes the local coefficient
is phase-independent. The redistribution experiment measured the worst cell
getting \emph{eight times worse}, which establishes the opposite. The $0.233$~kg
is therefore an artifact of the two particular grids compared, and two points
cannot distinguish that model from $E\propto(\sum h^3)^p$ or from a model with a
nonzero asymptote---both of which extrapolate very differently.

\paragraph{Switch count is not branch identity.} The comparison rested on
$19\to19$ switches. This document itself records a row that held $75$ switches
\emph{across} a genuine branch change; the constraint was stated in one section
and disregarded in another.

\paragraph{A state residual is not an objective error.} The $R_x/h^3$
constancy was offered as independent support for the objective-error model. It
is a scaling of the \emph{state} residual, and the only bridge to objective
error is the first-order sensitivity sum---which had already failed by a factor
of $5.3$.

\paragraph{That failure is worse than miscalibration.} The reviewers' reading,
which we accept: with bang-bang control on a mesh whose breakpoints are not
switch-aligned, the objective is \emph{non-smooth} in node placement, so a
smooth adjoint sensitivity does not compute a poorly-scaled gradient---it
computes a fictitious one. $\sum_k\lambda_k^{\!\top}d_k$ is not applicable here
at all.

\paragraph{Consequently withdrawn:} the $0.233$~kg point estimate; the
``basin dominates by $4.4\times$'' ratio, which loses its denominator; and
``the production mesh is good enough,'' which was never given an acceptance
tolerance against which ``enough'' could be judged.

\paragraph{One finding was challenged and survived.} A reviewer objected that
the median dilutes sparse switch spikes into a large population of smooth
cells. Re-tested on sums, switch cells carry $16.2\%$, $15.3\%$ and $10.5\%$ of
the total residual against population shares of $9.8\%$, $10.7\%$ and $9.8\%$:
mildly over-represented, not dominant, and the over-representation \emph{shrinks}
as switches multiply. The objection was fair; Finding~3.3 stands.

\paragraph{And one claim of ours was correctly demolished.} We had inferred
that basin fragility increases with switch count, because the two deeper rows
voided. Both reviewers objected that the same perturbation is a far larger
shock to a $171$-switch mesh than to a $19$-switch one. Testing at a $15\%$
blend: a $2\%$ change in $\sum h^3$ still flipped $76\to75$. The claim is
withdrawn, and what replaces it is sharper---these rows are \emph{knife-edge},
and essentially any mesh change reorganizes them.

\subsection{The recurring error, which is the transferable lesson}

This study has now made the same class of mistake twice.

\begin{enumerate}
  \item Phase~2 fitted Richardson orders across a branch change, when the
  study's own policy forbade exactly that.
  \item Phase~3 extrapolated $E = k\sum h^3$ assuming a phase-independent
  coefficient, when the study's own measurement had already falsified
  phase-independence.
\end{enumerate}

In both cases the contradicting evidence was \emph{already in hand and already
written down in the same document}. The failure was never missing data; it was
failing to apply our own stated constraint to our own new claim. The Phase~2
remedy was to move the rule out of a comment and into the arithmetic, where it
could not be forgotten. The analogous remedy here is a check that refuses to
report an extrapolated error whenever the measured spread of the local
coefficient exceeds a threshold---which, on this data, it plainly does.

\section{Experiments still to run}
\label{sec:next}

\begin{enumerate}
  \item \textbf{A branch-identity test that is not switch count.} Everything
  downstream needs this: matched switch \emph{times}, arc correspondence, and
  state/costate proximity between two meshes. Without it no same-branch
  comparison can be claimed, and both retractions in this document trace back
  to its absence.
  \item \textbf{Three or more same-branch meshes} spanning a range of
  $\sum h^3$, with competing error models fitted and compared rather than one
  assumed. Two points can support a bound; they cannot support a model.
  \item \textbf{A fixed-structure multi-phase formulation}, with switch times as
  explicit variables. This makes branch identity structural rather than
  detected, and it simultaneously supplies the switch-aligned mesh whose order
  the study has been trying to measure indirectly throughout.
  \item \textbf{Cislunar adapters} for the tulip and ELFO campaigns, each of
  which needs its own: the tulip builds its dynamics inline with no reusable
  right-hand side and stores no duals, and ELFO uses a free-$t_f$ two-primary
  nine-state solver.
\end{enumerate}

\noindent The basin thread is not on this list because it is no longer a
measurement problem. Every certified row examined is beaten on its own mesh by a
different starting guess, a one-node perturbation suffices to find the better
one, and nothing in the certification path looks. That is an operational change,
and on the evidence assembled here it is worth more than another digit on the
mesh.

\section*{Reproducing this}
\addcontentsline{toc}{section}{Reproducing this}

\begin{longtable}{ll}
\toprule
Component & Path (under \texttt{orbit\_transfer/}) \\
\midrule
\endhead
Sub-grid switch location   & \texttt{verify\_common/mesh/mesh\_switch\_times.m} \\
Cross-mesh switch matching & \texttt{verify\_common/mesh/mesh\_match\_switches.m} \\
Order and Richardson       & \texttt{verify\_common/mesh/mesh\_order.m} \\
Mesh ladder (Earth)        & \texttt{verify\_common/mesh/mesh\_ladder\_mee.m} \\
Collection and report      & \texttt{verify\_common/mesh/mesh\_collect.m}, \texttt{mesh\_report.m} \\
Tier A driver              & \texttt{verify\_common/mesh/run\_mesh\_study\_mee.m} \\
Basin probe (Earth)        & \texttt{verify\_common/mesh/mesh\_basin\_probe.m} \\
Basin probe (CR3BP)        & \texttt{verify\_common/mesh/mesh\_basin\_probe\_cr3bp.m} \\
Running record             & \texttt{verify\_common/doc/mesh\_study\_tierA\_results.md} \\
Plan                       & \texttt{docs/superpowers/plans/2026-07-25-mesh-convergence-study.md} \\
\bottomrule
\end{longtable}

\begin{thebibliography}{9}
\bibitem{roache}
P.~J. Roache, \emph{Verification and Validation in Computational Science and
Engineering}, Hermosa, 1998. (Observed order of accuracy and the grid
convergence index.)

\bibitem{hager}
W.~W. Hager, ``Runge--Kutta Methods in Optimal Control and the Transformed
Adjoint System,'' \emph{Numerische Mathematik} \textbf{87} (2000) 247--282.

\bibitem{dhv}
A.~L. Dontchev, W.~W. Hager and V.~M. Veliov, ``Second-Order Runge--Kutta
Approximations in Control Constrained Optimal Control,'' \emph{SIAM Journal on
Numerical Analysis} \textbf{38} (2000) 202--226.

\bibitem{agamawi}
Y.~M. Agamawi, W.~W. Hager and A.~V. Rao, ``Mesh Refinement Method for Solving
Bang-Bang Optimal Control Problems Using Direct Collocation,''
arXiv:1905.11895.

\bibitem{bertrand}
R.~Bertrand and R.~\'Ep\'enoy, ``New Smoothing Techniques for Solving
Bang--Bang Optimal Control Problems,'' \emph{Optimal Control Applications and
Methods} \textbf{23} (2002) 171--197.

\bibitem{betts}
J.~T. Betts, \emph{Practical Methods for Optimal Control Using Nonlinear
Programming}, 2nd ed., SIAM, 2010.
\end{thebibliography}

\end{document}
